\documentclass[12pt,a4paper]{article}

% Codificacao, idioma e tipografia
\usepackage[T1]{fontenc}
\usepackage[utf8]{inputenc}
\usepackage[brazilian]{babel}
\usepackage{lmodern}
\usepackage{microtype}

% Pagina, tabelas, figuras e equacoes
\usepackage{geometry}
\usepackage{graphicx}
\usepackage{booktabs}
\usepackage{amsmath}
\usepackage{enumitem}

% Citacoes e links
\usepackage[authoryear,round]{natbib}
\usepackage{hyperref}

\geometry{margin=2.5cm}
\graphicspath{{figuras/}}
\hypersetup{
  colorlinks=true,
  linkcolor=blue,
  urlcolor=blue,
  citecolor=blue
}

\title{Título do artigo}
\author{Nome do aluno\\Instituição}
\date{}

\begin{document}

\maketitle

\begin{abstract}
Apresente, em um único parágrafo, o contexto, a lacuna, o objetivo, o método,
o principal resultado e a conclusão. Durante as primeiras entregas, mantenha
este texto como um rascunho e revise-o após concluir as demais seções.
\end{abstract}

\noindent\textbf{Palavras-chave:} mobilidade aérea urbana; palavra 2;
palavra 3; palavra 4; palavra 5.

\section{Introdução}
\label{sec:introducao}

% Sugestão de sequência: contexto -> problema -> evidência -> lacuna -> pergunta.
Contextualize a transformação ou o fenômeno investigado e delimite seu recorte.
Explique por que o problema importa para a mobilidade aérea urbana e sustente a
motivação com fontes verificáveis. Por exemplo, uma revisão pode ser citada
entre parênteses \citep{garrow2021uam}.

Apresente o que a literatura já explica e identifique precisamente o que ainda
não foi respondido. Encerre a seção com o objetivo e a pergunta de pesquisa.

\begin{quote}
\textbf{Pergunta de pesquisa:} como [fator] influencia [resultado] no contexto
de [recorte]?
\end{quote}

% Remova este comentário depois de verificar:
% - o recorte está explícito;
% - a lacuna decorre das referências apresentadas;
% - a pergunta pode ser respondida pelo método proposto.

\section{Revisão de literatura}
\label{sec:revisao}

Organize os trabalhos por conceitos, abordagens ou resultados comparáveis, em
vez de apresentar uma lista de resumos. Para cada conjunto de estudos, registre:

\begin{itemize}[leftmargin=*]
  \item qual pergunta ou problema foi investigado;
  \item quais dados e métodos foram utilizados;
  \item quais resultados e limitações foram encontrados;
  \item como essas evidências sustentam a lacuna do presente artigo.
\end{itemize}

Conclua a seção posicionando o artigo diante da literatura: indique o que será
mantido, comparado, adaptado ou testado de maneira diferente.

\section{Método}
\label{sec:metodo}

Descreva o desenho da pesquisa de modo que outra pessoa possa compreender e,
quando aplicável, reproduzir a análise. Informe as fontes de dados, hipóteses,
variáveis, critérios de seleção, ferramentas, procedimentos e limitações.

Uma organização possível é:

\begin{enumerate}[leftmargin=*]
  \item definir o objeto e a unidade de análise;
  \item selecionar ou coletar os dados;
  \item aplicar o procedimento de análise;
  \item verificar a sensibilidade e interpretar os resultados.
\end{enumerate}

Se utilizar uma equação, defina cada símbolo logo após apresentá-la. Exemplo:

\begin{equation}
  C = \frac{N}{T},
  \label{eq:capacidade}
\end{equation}

em que $C$ representa a capacidade observada, $N$ o número de operações e $T$
o intervalo de tempo analisado. Substitua este exemplo pelo modelo do artigo.

\section{Resultados e discussão}
\label{sec:resultados}

Apresente os resultados na mesma ordem das etapas do método. Antes de cada
tabela ou figura, explique o que o leitor deve observar. Depois, interprete o
resultado, compare-o com a literatura e discuta implicações e limitações.

% Exemplo de figura. Coloque o arquivo em figuras/ e remova o símbolo %.
% \begin{figure}[htbp]
%   \centering
%   \includegraphics[width=0.78\textwidth]{nome-da-figura.png}
%   \caption{Título informativo da figura.}
%   \label{fig:resultado-principal}
% \end{figure}

% Exemplo de tabela editável.
\begin{table}[htbp]
  \centering
  \caption{Exemplo de comparação entre cenários.}
  \label{tab:cenarios}
  \begin{tabular}{lcc}
    \toprule
    Cenário & Indicador 1 & Indicador 2 \\
    \midrule
    Referência & -- & -- \\
    Alternativo & -- & -- \\
    \bottomrule
  \end{tabular}
\end{table}

Substitua a tabela de exemplo e evite repetir no texto todos os valores que já
estão visíveis. Destaque padrões, diferenças e respostas à pergunta de pesquisa.

\section{Conclusão}
\label{sec:conclusao}

Retome a pergunta de pesquisa e responda-a diretamente com base nos resultados.
Sintetize as contribuições teóricas e práticas, reconheça as limitações do
estudo e indique próximos passos coerentes com o que foi demonstrado.

Não introduza resultados ou referências essenciais apenas nesta seção. A
conclusão deve fechar o argumento construído no restante do artigo.


\bibliographystyle{plainnat}
\bibliography{referencias/referencias}

\end{document}
